% SORT -- os/samir/cmd/mutate.c d_sort (TO <file> ON <field>, default .DBF;
% ascending; Character, Numeric or Date field; a memo field refuses, 34;
% unknown field 48; character key over 100, 23; SAFETY ON refuses an
% existing target, 7; records taken in master-index order, equal keys kept
% in that order).  Lane samirdml 14a8fc75.  Gates: test-samir-dml,
% test-samir-dml-emu.
\begin{XbaseEntry}{SORT}{command}
Copies the database in USE to a new database, in the order of a field.
\Syntax
\SyntaxLine{\kw{SORT TO} \xp{file} \kw{ON} \xp{field}}
\Usage
\cmd{SORT} writes every record of the database in the selected work area to
a new database, in ascending order of the field named.  If no extension is
given, \file{.DBF} is assumed.  The field must be of type Character, Numeric
or Date; records with equal values keep the order in which they were read.
The database in USE is not changed.

A database with a memo field cannot be sorted (\xmsg{34}{Operation with
Memo field invalid.}).  A field the database does not have is reported as
\xmsg{48}{Field not found.}, and a character field of more than 100
characters as \xmsg{23}{Index is too big (100 char maximum).}  While SET
SAFETY is ON, a database of the new name is not replaced
(\xmsg{7}{File already exists.}).
\Examples
The ledger in order of due date.  The due dates keyed with the year 00 are
stored as 1900 (see \cmd{SET CENTURY}), and so come first:
\DotExample{examples/c-sort}
\file{STAFF.DBF} has a memo field, \texttt{NOTES}, and cannot be sorted:
\DotExample{examples/c-sort-2}
\SeeAlso{INDEX, TOTAL}
\end{XbaseEntry}

